\documentclass[11pt,a4paper]{article}
\usepackage[T1]{fontenc}
\usepackage{lmodern}
\usepackage{amsmath,amssymb,amsthm}
\usepackage{microtype}
\usepackage[margin=24mm]{geometry}
\usepackage{booktabs}
\usepackage[hidelinks]{hyperref}
\hypersetup{pdftitle={An elementary obstruction for y cubed plus xy equals x to the fourth plus x plus 3},pdfauthor={},pdfsubject={A self-contained proof that the equation has no integer solutions}}
\newtheorem{lemma}{Lemma}
\newtheorem{theorem}{Theorem}
\newcommand{\Z}{\mathbb Z}
\newcommand{\Fp}{\mathbb F_p}
\setlength{\parindent}{0pt}
\setlength{\parskip}{4pt}
\setlength{\emergencystretch}{2em}
\title{\vspace{-1.1cm}An elementary obstruction for\\[3pt]
\(y^3+xy=x^4+x+3\)}
\author{}
\date{\vspace{-1.1cm}}

\begin{document}
\maketitle
\vspace{-0.4cm}

The equation has \textbf{no integer solutions}. The obstruction comes from two
norm factorizations: they force two positive odd integers to be congruent to
$1$ or $3$ modulo $8$, although their product is congruent to $7$ modulo $8$.

\section{An elementary prime-divisor lemma}

\begin{lemma}\label{lem:prime}
If $p$ is a prime with $p\equiv5$ or $7\pmod8$ and
$p\mid a^2+2b^2$, then $p\mid a$ and $p\mid b$.
Consequently, if a positive odd integer has no prime divisor congruent to
$5$ or $7$ modulo $8$, it is congruent to $1$ or $3$ modulo $8$.
\end{lemma}

\begin{proof}
We first show that $r^2=-2$ in $\Fp$ implies $p\equiv1$ or $3\pmod8$.
Choose a root $u$ of $T^2-rT-1$ in a field containing $\Fp$.
Such a field is either $\Fp$ itself or the quotient by this quadratic if it
is irreducible. Then
\[
 u^2=ru+1,\qquad u^3=r-u,\qquad u^4=-1.
\]
Thus $u$ has multiplicative order exactly $8$. The two roots of the
quadratic are $u$ and $u^3$. In characteristic $p$, the binomial theorem
and $r^p=r$ give
\[
 (u^p)^2-r u^p-1=(u^2-ru-1)^p=0.
\]
Here $r^p=r$ is elementary: multiplication by any nonzero residue
permutes the nonzero residues, and multiplying them proves
$r^{p-1}=1$. Hence $u^p=u$ or $u^p=u^3$, so $p\equiv1$ or $3\pmod8$.

Now, if $p\nmid b$ and $a^2+2b^2\equiv0\pmod p$, then
$r=ab^{-1}$ would satisfy $r^2=-2$, a contradiction. Thus $p\mid b$,
and then $p\mid a$. Finally, a positive odd integer with no forbidden
prime divisor is a product of primes congruent to $1$ or $3$ modulo $8$;
such a product is again $1$ or $3$, since $3^2\equiv1\pmod8$.
\end{proof}

\section{Nonexistence theorem}

\begin{theorem}
There is no pair $(x,y)\in\Z^2$ satisfying
\begin{equation}\label{eq:main}
 y^3+xy=x^4+x+3.
\end{equation}
\end{theorem}

\begin{proof}
Suppose that \eqref{eq:main} holds.

\textit{Positivity and parity.}
For every integer $x$, one has $x^4+x+3\geq3$.
If $y\leq0$, positivity of the left side would force
$x=-u$, $y=-v$ with $u,v>0$ and $u>v^2$. In particular $v\leq u-1$, giving
\[
 y^3+xy=v(u-v^2)\leq u(u-1)<u^4-u+3=x^4+x+3,
\]
where the strict inequality follows from $u^4-u^2+3>0$.
This is impossible. Also $y=1$ would give $x^4=-2$. Thus $y>1$.

Modulo $2$, equation \eqref{eq:main} gives $y(1+x)\equiv1$, so
$x$ is even and $y$ is odd. Modulo $8$, therefore,
\[
 (1+x)y\equiv x+3\pmod8.
\]
Every odd residue is its own inverse modulo $8$, and for even $x$ one has
$x^2\equiv2x\pmod8$. It follows that
\begin{equation}\label{eq:parity}
 y\equiv(1+x)(x+3)\equiv2x+3\pmod8.
\end{equation}

\textit{The first factor.}
Set
\[
 D=y^2+y+1+x.
\]
Equation \eqref{eq:main} gives
\begin{equation}\label{eq:first}
 (y-1)D=x^4+2=(x^2)^2+2\cdot1^2.
\end{equation}
Since $y>1$, we have $D>0$; it is also odd. If a prime
$p\equiv5$ or $7\pmod8$ divided $D$, Lemma~\ref{lem:prime} applied to
$x^4+2$ would imply $p\mid1$, a contradiction. Consequently,
\begin{equation}\label{eq:Dres}
 D\equiv1\ \text{or}\ 3\pmod8.
\end{equation}

\textit{The second factor.}
Define
\[
 \ell=2x-13,\qquad
 C=y^2+y\ell+\ell^2+x,\qquad
 A=x^2-4x+20,\qquad B=7x-30.
\]
The key identity is
\begin{equation}\label{eq:second}
 (y-\ell)C
 =x^4+x+3-\ell^3-x\ell
 =A^2+2B^2.
\end{equation}
Moreover, $C$ is odd, and it is positive because
\begin{equation}\label{eq:positive}
 12C=3(2y+2x-13)^2+4(3x-19)^2+77>0.
\end{equation}

We claim that $C$ has no prime divisor $p\equiv5$ or $7\pmod8$.
Indeed, such a divisor would divide $A^2+2B^2$, and
Lemma~\ref{lem:prime} would force $p\mid A$ and $p\mid B$.
But
\begin{equation}\label{eq:exceptional}
 49A=B(7x+2)+1040,\qquad 1040=2^4\cdot5\cdot13,
\end{equation}
so $p=5$ or $p=13$.

For $p=5$, the condition $p\mid B$ gives $x\equiv0\pmod5$, hence
$\ell\equiv2\pmod5$. Therefore $p\mid C$ would imply
\[
 0\equiv C\equiv y^2+2y+4=(y+1)^2+3\pmod5,
\]
so $(y+1)^2\equiv2\pmod5$, impossible since the square residues are
$0,1,4$.

For $p=13$, the condition $p\mid B$ gives $x\equiv8\pmod{13}$,
hence $\ell\equiv3\pmod{13}$. Thus $p\mid C$ would imply
\[
 0\equiv C\equiv y^2+3y+4\pmod{13},
 \qquad (2y+3)^2\equiv6\pmod{13}.
\]
This too is impossible: the square residues modulo $13$ are
$0,1,3,4,9,10,12$. The claim is proved, and Lemma~\ref{lem:prime} now gives
\begin{equation}\label{eq:Cres}
 C\equiv1\ \text{or}\ 3\pmod8.
\end{equation}

\textit{The contradiction.}
Since $\ell\equiv2x+3\equiv y\pmod8$ and $y$ is odd,
\[
 D\equiv3x+5\pmod8,\qquad C\equiv x+3\pmod8.
\]
The four possible residues of the even integer $x$ give
\[
\begin{array}{c|rrrr}
 x\pmod8&0&2&4&6\\ \hline
 D\pmod8&5&3&1&7\\
 C\pmod8&3&5&7&1
\end{array}
\]
No column is consistent with both \eqref{eq:Dres} and \eqref{eq:Cres}.
Equivalently, $DC\equiv7\pmod8$, whereas a product of two residues from
$\{1,3\}$ must lie in $\{1,3\}$. This contradiction proves the theorem.
\end{proof}

\section{Motivation for the construction}

The useful starting point is not to compare the sizes of $y^3$ and $x^4$,
but to exploit the shared term $xy-x=x(y-1)$. Subtracting the value at
$y=1$ produces \eqref{eq:first}. The expression $x^4+2$ is of the form
$a^2+2b^2$ with $b=1$, so it cannot have a prime divisor congruent to $5$
or $7$ modulo $8$. This forces the first odd factor $D$ into the
residues $1,3$ and, by the last table, restricts $x$ to $2,4$ modulo $8$.

A second factorization must exclude precisely those remaining residues.
Equation \eqref{eq:parity} suggests subtracting a linear value
$\ell=2x+b$ with $b\equiv3\pmod8$. Its complementary factor
$y^2+y\ell+\ell^2+x$ is then automatically $x+3$ modulo $8$, exactly the
needed opposite restriction.

To construct a norm identity, seek integers $u,v,w,b$ such that
\[
 x^4+x+3-(2x+b)^3-x(2x+b)
 =(x^2-4x+u)^2+2(vx+w)^2.
\]
The coefficient $-4$ is forced by the cubic term. Matching the remaining
coefficients gives
\[
 \begin{aligned}
 u+v^2&=-6b-9,\\
 -8u+4vw&=1-b-6b^2,\\
 u^2+2w^2&=3-b^3.
 \end{aligned}
\]
The tuple $(b,u,v,w)=(-13,20,7,-30)$ satisfies these three identities
and has the required $b\equiv3\pmod8$. It therefore produces
\[
 x^4+x+3-(2x-13)^3-x(2x-13)
 =(x^2-4x+20)^2+2(7x-30)^2.
\]
The two norm coordinates can share a forbidden odd prime only at $5$ or
$13$, as \eqref{eq:exceptional} shows, and both exceptions are eliminated
by the displayed square-residue checks. Thus the second factor also has
only the allowed prime divisors. The decisive step is the combination
of these two norm restrictions, not a congruence test on the original
equation alone.

\end{document}
